%%%%%%%%%%%%%%%%%%%%%%%%%%%%%%%%%%%%%%%%%%%%%%%%%%%%%%%%%%%%%%%%%%%%%%%%%%%%%%%%%%%

\blankpage
\pagebreak
\thispagestyle{empty}
\addcontentsline{toc}{part}{Fès}
\mbox{}
\begin{center}

\vspace{3cm}

{\brabo\HUGE\itshape Fès}
\end{center}
\pagebreak

%%%%%%%%%%%%%%%%%%%%%%%%%%%%%%%%%%%%%%%%%%%%%%%%%%%%%%%%%%%%%%%%%%%%%%%%%%%%%%%%%%%

\chapter{Bab Boujloud}

Bab Boujloud é uma das principais portas de entrada da medina de Fès. A porta atual foi construída no início do século \textsc{xx}, substituindo uma entrada mais antiga, mas mantém a linguagem decorativa tradicional, com azulejos e motivos geométricos. A partir daqui começa a Talaa Kebira, uma das principais ruas da medina: ruas estreitas, comércio, oficinas, mesquitas, escolas religiosas, fontes e casas tradicionais se sucedem em um tecido urbano em que é difícil separar o espaço comercial do residencial.

\chapter{Talaa Kebira e Talaa Sghira}

A Talaa Kebira é a mais conhecida e passa por alguns dos monumentos mais importantes da cidade, enquanto a Talaa Sghira oferece uma alternativa menos movimentada. Caminhar por essas ruas é parte essencial da visita. Em vez de tentar seguir apenas uma lista de monumentos, vale observar como os espaços mudam à medida que se passa dos grandes eixos comerciais para pequenas ruas residenciais, áreas de artesãos, mesquitas e \textit{funduqs}.

\chapter{Madrassa Bou Inania}

A Madrassa Bou Inania foi construída entre 1351 e 1356, durante a dinastia merínida, e é uma das obras mais importantes da arquitetura religiosa de Fès. Além de funcionar como escola religiosa, a madrassa fazia parte da infraestrutura intelectual da cidade. O edifício organiza-se em torno de um pátio central e reúne alguns dos elementos mais característicos da arquitetura marroquina: \textit{zellij}, estuque esculpido, madeira de cedro e inscrições caligráficas.

\chapter{Dar al-Magana}

Dar al-Magana está associada ao antigo relógio hidráulico de Fès. O sistema medieval utilizava recipientes e mecanismos de água para marcar a passagem das horas. A estrutura que sobrevive hoje está bastante modificada, mas o local é interessante para compreender a sofisticação técnica da cidade medieval e a relação entre ciência, religião e organização da vida cotidiana.

\chapter{Museu Nejjarine}

O Museu Nejjarine de Artes e Ofícios da Madeira funciona em um antigo \textit{funduq}, uma hospedaria destinada a comerciantes e viajantes. O edifício é interessante tanto pela coleção quanto pela arquitetura. A exposição reúne objetos relacionados ao trabalho da madeira, enquanto o próprio funduq permite imaginar a circulação de mercadores e artesãos dentro da medina. O terraço também oferece uma boa perspectiva sobre a medina.

\chapter{Fonte Nejjarine}

A fonte de Nejjarine é uma das fontes decoradas mais conhecidas de Fès. O conjunto combina água, madeira e \textit{zellij}, mostrando como elementos utilitários podiam receber tratamento arquitetônico elaborado. Fontes públicas faziam parte da infraestrutura cotidiana da medina e também funcionavam como elementos de composição urbana.

\chapter{Al-Qarawiyyin}

A Mesquita e Universidade de Al-Qarawiyyin é um dos símbolos intelectuais de Fès e um dos grandes centros históricos de conhecimento do mundo islâmico. Fundada no século \textsc{ix} por Fatima al-Fihri, tornou-se um importante centro de estudos religiosos e jurídicos. Fès cresceu em torno de instituições desse tipo, que ajudaram a estabelecer sua reputação como centro intelectual do Magrebe. A entrada no espaço religioso é restrita a não muçulmanos, portanto a visita costuma ser feita a partir do exterior.

\chapter{Madrassa Al-Attarine}

A Madrassa Al-Attarine foi construída no século \textsc{xiv}, durante o período merínida, e é uma das obras mais refinadas da arquitetura de Fès. Como na Bou Inania, o edifício combina \textit{zellij}, estuque, madeira de cedro e inscrições caligráficas. A escala mais compacta permite observar de perto os detalhes da ornamentação.

\chapter{Praça Seffarine}

A Praça Seffarine fica junto ao rio e é conhecida pelas oficinas de artesãos que trabalham o cobre. O interesse da praça está em observar o trabalho artesanal. O som dos martelos, os objetos sendo moldados e a presença das oficinas ajudam a entender como o artesanato continua integrado à vida cotidiana da medina.

\chapter{Souks e bairros de artesãos}

Os souks de Fès são organizados em diferentes áreas de especialização. Vale observar os bairros onde os artesãos continuam trabalhando. Fès é especialmente conhecida por sua cerâmica azul e por uma longa tradição de trabalhos em madeira, metal, couro e tecidos.

\begin{itemize}
\item Praça Seffarine: cobre e metal;
\item Nejjarine: madeira;
\item áreas de cerâmica: produção do tradicional azul de Fès;
\item mercados de tecidos e bordados;
\item oficinas de couro;
\item perfumarias e especiarias.
\end{itemize}

\chapter{Fès el-Jdid}

Fès el-Jdid, ou ``Fès Nova'', foi criada pelos merínidas no século \textsc{xiii}, quando a cidade passou a receber novas estruturas palacianas e administrativas. Enquanto a medina antiga concentra comércio, instituições religiosas, artesãos e vida residencial, Fès el-Jdid foi construída em torno do poder político e do complexo do palácio real.

\chapter{Mellah de Fès}

O bairro judaico foi estabelecido em Fès el-Jdid. Sua localização junto ao complexo real está relacionada à organização política e social do Marrocos islâmico, na qual as comunidades judaicas tinham uma relação particular com o sultão. As casas judaicas podiam ter varandas voltadas para a rua, enquanto muitas casas da medina eram mais fechadas em relação ao espaço público. 

Há também um importante cemitério e a sinagoga Ibn Danan, uma das principais sinagogas históricas de Fès. Foi restaurada e conserva elementos do espaço religioso tradicional, incluindo a área de oração, a galeria feminina e objetos relacionados à vida comunitária.

\chapter*{Palácio Real\\ \& Bab el-Makhzen}
\addcontentsline{toc}{chapter}{Palácio Real \& Bab el-Makhzen}

O Palácio Real de Fès ocupa uma grande área de Fès el-Jdid e continua sendo uma residência real, portanto não é aberto normalmente à visitação. As portas de Bab el-Makhzen, entretanto, podem ser observadas do lado de fora e estão entre os elementos arquitetônicos mais conhecidos da cidade.

\chapter{Dar Batha}

O Dar Batha é um antigo palácio transformado em museu de artes tradicionais marroquinas. A coleção reúne objetos de cerâmica, madeira e outras artes decorativas. O próprio edifício é parte importante da visita, com pátios, jardins e elementos da arquitetura palaciana.

\chapter{Jnan Sbil}

Jnan Sbil é um grande jardim histórico criado no século \textsc{xix}, situado entre Fès el-Bali e Fès el-Jdid. Depois de passar várias horas na medina, é um dos melhores lugares para diminuir o ritmo. O jardim reúne caminhos, fontes, canais de água e vegetação abundante.

\chapter{Túmulos Merínidas}

Os Túmulos Merínidas ficam em uma colina acima de Fès e oferecem uma das melhores vistas panorâmicas da cidade. As estruturas funerárias estão em ruínas, mas o principal interesse da visita é a perspectiva sobre a medina: as casas, minaretes, muralhas e bairros se espalham pela paisagem abaixo.

\chapter{Outros pontos importantes}

A cerâmica é uma das tradições artesanais mais conhecidas da cidade. Fès é especialmente associada ao azul profundo utilizado em pratos, azulejos e objetos decorativos. 

A cidade também conserva uma longa tradição de produção e transmissão de textos religiosos, jurídicos e literários. Vale procurar por exposições temporárias, bibliotecas, centros culturais e coleções de manuscritos que estejam abertos ao público nas datas da visita.

Por, fim, Fès possui uma tradição importante de música arabo-andalusina, uma das manifestações culturais que melhor permitem pensar na continuidade entre o Magrebe e Al-Andalus.